
\begin{examplebox}{Example}
As a practical demonstration of our recommendations, we discuss the GSM8K benchmark \cite{cobbe_training_2021}. We chose this benchmark because of its widespread adoption and examples of simple ways to address many, but not all, of our recommendations.\\

\textbf{Define the phenomenon:} 
GSM8K describes itself as `grade school math problems...using basic arithmetic' which are `useful for probing the informal reasoning ability
 of large language models'. This definition describes both the phenomenon, [mathematical] reasoning, and the task, arithmetic-based math problems, making the operationalisation of `reasoning' and the scope of the assessment clear. The authors also note that the difficulty comes from a combination of `reading comprehension' and `logical reasoning'. Based on our checklist, we would recommend assessing the relative impact of each of these skills. One approach would be to annotate the relative demand for each skill on each of the problems, similar to \citeauthor{zhou2025general} \cite{zhou2025general}.\\

\textbf{Measure only the phenomenon:} 
Beyond reading comprehension and logical reasoning, GSM8K requires arithmetic calculations. The authors note this as a potential confounder and provide a calculator tool to their model. However, we would recommend building this into the benchmark directly to avoid score differences due to testing setups \cite{bidermanLessonsTrenchesReproducible2024a}. Instead, the answer could be provided as a mathematical expression which is evaluated as part of the scoring. The answer format itself is simple, although the potential impacts of tokenization should be considered \cite{singh2024tokenization}.\\

\textbf{Construct a representative dataset for the task:} 
Annotators creating the dataset are given a diverse set of seed prompts and the dataset is tested for reliability by human annotators prior to release. Additional tests to target known LLM weaknesses, such as re-wording the questions, would improve the robustness of the results~\cite{mirzadehGSMSymbolicUnderstandingLimitations2024}.\\

\textbf{Acknowledge limitations of reusing datasets:} 
The dataset is created from scratch for this project, which avoids the issues of reused datasets.\\

\textbf{Prepare for contamination:} 
We recommend the addition of a canary string and a set of held-out task items. Testing performance on a new set of similar benchmark questions also shows significant performance drops, likely indicating that contamination has occurred \cite{zhang2024careful}.\\

\textbf{Use statistical methods to compare models:} 
The benchmark reports a large sample size justified by the deliberate creation of diverse questions. Comparisons between the models being tested are reported over multiple runs with standard deviations.\\

\textbf{Conduct an error analysis:} 
No error analysis is conducted. By creating a typology of error types on GSM8K, other works were able to guide improvements in future models \cite{wei2022chain, wang-etal-2023-plan}.\\

\textbf{Justify construct validity:} 
The authors describe how their dataset differs from existing sets of maths problems in quality and scale. We would also recommend a discussion of how maths reasoning problems relate to logical reasoning broadly given the stated requirements of `reading comprehension' and `logical reasoning'.\\

\textbf{Conclusion:}
GSM8K is a generally valid benchmark for measuring performance on grade school maths questions. Contamination has likely been a factor in increased scores over time which may have been partially avoidable. An error analysis would also have furthered the usefulness of the benchmark. Where interpretation stretches from `math problems' to `logical reasoning', greater clarity would be valuable to help readers interpret the results.

\end{examplebox}